\lyseis{Γ3.2}{Δενδροδιάγραμμα και πίνακας}
\sol{1} Από τη ρίζα ξεκινούν δύο κλαδιά (Κ και Γ) και από το καθένα δύο νέα (Κ και Γ). Τα μονοπάτια δίνουν $\{\mathrm{ΚΚ},\ \mathrm{ΚΓ},\ \mathrm{ΓΚ},\ \mathrm{ΓΓ}\}$, $4$ στοιχεία.
\sol{2} Η γραμμή $k$ και η στήλη $\ell$ έχουν άθροισμα $k+\ell$. Άθροισμα $8$: τα $(2,6),(3,5),(4,4),(5,3),(6,2)$, δηλαδή $5$ ζευγάρια. Διαφορετικά αθροίσματα: $2,3,\dots,12$, δηλαδή $11$.
\sol{3} α) $2\cdot2\cdot2=8$. β) $6\cdot2=12$. γ) $6\cdot6=36$. δ) $2^4=16$.
\sol{4} Μέγεθος: μικρή, μεγάλη. Από κάθε μέγεθος ξεκινούν τρία κλαδιά (μαργαρίτα, ζαμπόν, μανιτάρια). Επιλογές: μικρή-μαργαρίτα, μικρή-ζαμπόν, μικρή-μανιτάρια, μεγάλη-μαργαρίτα, μεγάλη-ζαμπόν, μεγάλη-μανιτάρια. Συνολικά $2\cdot3=6$.
\sol{5} α) Σωστό. β) Λάθος: $2\cdot2\cdot2=8$. γ) Σωστό. δ) Λάθος: $6\cdot6=36$. ε) Σωστό: $4\cdot5=20$.
\sol{6} $4\cdot3=12$ συνδυασμοί.
\sol{7} Από τη ρίζα ξεκινούν τρία κλαδιά (Μ, Κ, Λ) και από το καθένα δύο κλαδιά (οι άλλες δύο καραμέλες). Δειγματικός χώρος: $\{\mathrm{ΜΚ},\ \mathrm{ΜΛ},\ \mathrm{ΚΜ},\ \mathrm{ΚΛ},\ \mathrm{ΛΜ},\ \mathrm{ΛΚ}\}$, $3\cdot2=6$ στοιχεία.
\sol{8} Από τη ρίζα ξεκινούν $2$ κλαδιά (ψωμί), από το καθένα $3$ (γέμιση) και από κάθε γέμιση $2$ (ρόφημα). Πρωινά: $2\cdot3\cdot2=12$.
\sol{9} α) $10\cdot10\cdot10\cdot10=10\,000$. β) Το πρώτο ψηφίο είναι το $1$. Το δεύτερο έχει $4$ επιλογές ($2,3,4,5$) και το τρίτο $3$ (ό,τι έμεινε). Άρα $1\cdot4\cdot3=12$ κωδικοί.
\sol{10} Η διαφορά είναι $0$ όταν τα δύο ζάρια φέρουν το ίδιο: $(1,1),(2,2),\dots,(6,6)$, $6$ ζευγάρια. Διαφορά $5$: $(1,6)$ και $(6,1)$, δηλαδή $2$ ζευγάρια.
\sol{11} Το δενδροδιάγραμμα έχει $8$ φύλλα. Ακριβώς μία εύστοχη: $\mathrm{ΕΑΑ},\ \mathrm{ΑΕΑ},\ \mathrm{ΑΑΕ}$, $3$ αποτελέσματα. Τουλάχιστον δύο εύστοχες: $\mathrm{ΕΕΕ},\ \mathrm{ΕΕΑ},\ \mathrm{ΕΑΕ},\ \mathrm{ΑΕΕ}$, $4$ αποτελέσματα.
\sol{12} Πίνακας: $\nu=1$: $2$, $\nu=2$: $4$, $\nu=3$: $8$, $\nu=4$: $16$, $\nu=5$: $32$. Κάθε ρίψη διπλασιάζει τα αποτελέσματα, άρα το πλήθος είναι $2^\nu$. Για $10$ ρίψεις: $2^{10}=1024$.
\sol{13} Όταν μετράμε ζευγάρια (κόκκινο, μπλε), το $(3,4)$ και το $(4,3)$ είναι διαφορετικά, άρα δίνει δίκιο στον Ηλία. Όταν ενδιαφέρει μόνο το άθροισμα, τα δύο ζευγάρια δίνουν το ίδιο άθροισμα $7$, και ο Βασίλης έχει δίκιο μόνο ως προς αυτό.
\sol{14} Το πρώτο ψηφίο έχει $4$ επιλογές, το δεύτερο $3$, το τρίτο $2$. Άρα $4\cdot3\cdot2=24$ αριθμοί.
\sol{15} Λείπει το $\mathrm{ΓΚ}$. Στο δενδροδιάγραμμα από το $\mathrm{Γ}$ ξεκινούν δύο κλαδιά (Κ και Γ), όπως και από το Κ. Έτσι τα φύλλα είναι $4$ και όχι $3$.
\sol{16} Ελεύθερη εργασία. Π.χ. $3$ ζευγάρια παπούτσια, $2$ μπουφάν και $2$ σκούφοι δίνουν $3\cdot2\cdot2=12$ επιλογές. Το δενδροδιάγραμμα πρέπει να έχει τόσα φύλλα όσο το γινόμενο.
\sol{17} Ο πίνακας έχει $3\cdot3=9$ κελιά, άρα $9$ αποτελέσματα. Ισοπαλία έχουμε όταν οι δύο παίκτες διαλέγουν το ίδιο: πέτρα-πέτρα, ψαλίδι-ψαλίδι, χαρτί-χαρτί, δηλαδή σε $3$ από τα $9$. Στα άλλα $6$ κερδίζει είτε ο πρώτος (σε $3$) είτε ο δεύτερος (σε $3$).
